
%%% fancy intro/abstract
~\vfill{}
\newcommand\forcelink[2] {\href{#1}{{\color{blue!70!black}#2}}}
\begin{center}
	\parbox[c]{0.8\linewidth}{
	\large
	\justifying{}

	\RpgDropCap{W}{elcome to the \rpgtex{} package}. This \LaTeX{} package is designed to allow users to flexibly typeset documents associated with Role Playing Games such as \textit{Dungeons \& Dragons} -- and many more besides. This packages defines a central engine: \texttt{rpgcore} which define a number of useful functions and classes, and a flexible set of \texttt{themes} which control how those commands are rendered in the final document.


	\vspace{1cm}

	\subsection*{Getting Started}

		This document is the full reference manual. For a quick example of how to get started with your first \rpgtex{} document, see the \forcelink{https://github.com/DrFraserGovil/rpgtex/tree/main/example}{example documents}, which demonstrate the basic features of this package in a readable fashion.

		\vspace{1cm}



	\subsection*{Attribution \& License}

		This package would not have been possible without the team who developed \forcelink{https://github.com/rpgtex/DND-5e-LaTeX-Template/tree/dev}{its predecessor, the `DND 5e LateX Template'}. That code was released under an MIT license, the text of which can be found in the LICENSE file. \texttt{rpgtex} is released under the same license.
		\vspace{1cm}
	\subsection*{AI Usage Declaration}

		The author used a handful of AI chat tools during development to help guide their thinking. The package itself was written entirely by hand: no AI tools were used to write any of its code. An AI assistant (Claude) helped edit this documentation. All mistakes and bad decisions are organic in nature.
		}
\end{center}

\vfill{}
